\documentclass[11pt]{book}
\title{Multipage Book for Partition Benchmarks}
\author{Ferritex Bench Corpus}
\date{April 2026}

\newcommand{\bookparagraph}[2]{%
  \noindent\textbf{#1.} #2\par\medskip
}
\newcommand{\fillchapter}[2]{%
  \bookparagraph{#1}{#2}%
  \bookparagraph{#1}{#2}%
  \bookparagraph{#1}{#2}%
  \bookparagraph{#1}{#2}%
  \bookparagraph{#1}{#2}%
  \bookparagraph{#1}{#2}%
  \bookparagraph{#1}{#2}%
  \bookparagraph{#1}{#2}%
  \bookparagraph{#1}{#2}%
}

\begin{document}
\frontmatter
\maketitle
\tableofcontents

\mainmatter
\chapter{Editorial Origins}
\fillchapter{Origin note}{
The first stewards of the benchmark collection treated each document as a
small public artifact rather than a disposable test case. They wrote clear
titles, preserved the context for unusual formatting choices, and kept a
record of how reference PDFs were produced. That discipline made the archive
more useful because later contributors could understand the purpose of a
fixture without reverse engineering it from raw source alone.
}

\chapter{Growth of the Corpus}
\fillchapter{Growth note}{
As the corpus expanded, the maintainers learned that structure mattered as
much as quantity. New subsets were easier to review when each one expressed
a distinct idea: layout features in one area, navigation behavior in
another, and partition-oriented documents in a third. Long-form book
fixtures were especially helpful because they exposed page sequencing,
chapter starts, and cumulative pagination over a broader range.
}

\chapter{Publication Discipline}
\fillchapter{Publication note}{
Publication week followed a familiar script. One maintainer confirmed the
inventory, another generated reference PDFs, and a third checked that the
bookmarks, page counts, and provenance notes still matched the source tree.
This division of labor reduced mistakes and gave reviewers multiple ways to
verify the same outcome before the artifacts were archived.
}

\chapter{Stewardship}
\fillchapter{Stewardship note}{
Sustained stewardship depends on small repeated actions: naming files
carefully, removing stale auxiliary outputs, and recording the exact command
used to build each reference. None of those steps is glamorous, yet together
they are what make a benchmark trustworthy over time. A long book fixture is
an effective place to capture that theme because it unfolds across many
pages and chapters.
}
\end{document}
